\documentclass[11pt,a4paper]{article}
\usepackage[utf8]{inputenc}
\usepackage[T1]{fontenc}
\usepackage[spanish,es-minimal,es-noshorthands]{babel}
\AtBeginDocument{\renewcommand{\sin}{\operatorname{sen}}\renewcommand{\arcsin}{\operatorname{arcsen}}}
\usepackage{lmodern}
\usepackage[sfdefault,lining]{FiraSans}
\usepackage[a4paper,top=2.3cm,bottom=2.2cm,left=1.8cm,right=1.8cm,headheight=1.3cm,headsep=0.4cm,footskip=1cm]{geometry}
\usepackage{amsmath,amssymb}
\usepackage{icomma}
\usepackage{graphicx}
\usepackage[table]{xcolor}
\usepackage{tikz}
\usepackage{fontawesome5}
\usepackage{ccicons}
\usepackage{fancyhdr}
\usepackage{lastpage}
\usepackage{tasks}
\usepackage{needspace}
\usepackage[most]{tcolorbox}
\usepackage{array}
\usepackage[hidelinks]{hyperref}
\AtBeginDocument{\hypersetup{pdftitle={Soluciones de la ficha de la unidad 1: Números naturales, enteros y fraccionarios. 4º ESO},pdfsubject={Matemáticas A 4º ESO},pdfkeywords={números naturales, operaciones combinadas, potencias, raíces, números enteros, valor absoluto, fracciones, fracciones equivalentes, simplificación, amplificación, problemas con fracciones, propiedades de las potencias}}}

\definecolor{navy}{HTML}{1B2748}
\definecolor{accent}{HTML}{E8603A}
\definecolor{teal}{HTML}{1F9E90}
\definecolor{soft}{HTML}{EEF1F8}
\definecolor{muted}{HTML}{5D6685}

\setlength{\parindent}{0pt}
\setlength{\parskip}{3pt}
\settasks{label=\alph*), label-format=\bfseries\color{navy}, label-width=1.3em, item-indent=1.8em, column-sep=0.6em, after-item-skip=2pt, before-skip=3pt, after-skip=4pt}

% ---------- capçalera i peu ----------
\newcommand{\llicencia}{\href{https://creativecommons.org/licenses/by-nc-sa/4.0/deed.es}{\ccbyncsa}\ \ Licencia Creative Commons Atribución-NoComercial-CompartirIgual 4.0 Internacional (CC BY-NC-SA 4.0)}
\pagestyle{fancy}
\fancyhf{}
\fancyhead[L]{\raisebox{-0.35\height}{\includegraphics[height=0.95cm]{logo-gva.png}}}
\fancyhead[C]{\small\color{muted}Soluciones de la ficha · Unidad 1: Números naturales, enteros y fraccionarios · 4\textsuperscript{o} ESO A}
\fancyhead[R]{\raisebox{-0.35\height}{\includegraphics[height=0.95cm]{logo-vermellar.png}}}
\fancyfoot[L]{\scriptsize\color{muted}\llicencia}
\fancyfoot[R]{\small\color{muted}\thepage}
\renewcommand{\headrulewidth}{0.4pt}
\renewcommand{\headrule}{\color{muted!40}\hrule height \headrulewidth}
\fancypagestyle{primera}{\fancyhf{}\renewcommand{\headrulewidth}{0pt}
  \fancyfoot[L]{\scriptsize\color{muted}\llicencia}
  \fancyfoot[R]{\small\color{muted}\thepage}}

% ---------- nivells i exercicis ----------
\newcommand{\estrelles}[1]{{\color{accent}\ifcase#1\or\faStar\faStar[regular]\faStar[regular]\or\faStar\faStar\faStar[regular]\or\faStar\faStar\faStar\fi}}
\newcounter{exer}
\newcommand{\ex}[1]{\par\needspace{4\baselineskip}\medskip\stepcounter{exer}\noindent{\color{navy}\bfseries\theexer.}\ \ {\small\estrelles{#1}}\ \ }

\newcommand{\seccio}[2]{\par\needspace{10\baselineskip}\bigskip
  \begin{tcolorbox}[colback=navy,colframe=navy,coltext=white,arc=2mm,boxrule=0pt,left=3mm,right=3mm,top=1.2mm,bottom=1.2mm]
  {\large\bfseries #1\quad #2}
  \end{tcolorbox}\vspace{-2pt}}

\newcolumntype{C}[1]{>{\centering\arraybackslash}p{#1}}

% quadre de recordatori (només a la fitxa d'exercicis)
\newcommand{\recorda}[1]{\par\smallskip\begin{tcolorbox}[colback=accent!7,colframe=accent!55,arc=2mm,boxrule=0.5pt,left=2.5mm,right=2.5mm,top=1mm,bottom=1mm]\small{\color{accent}\faLightbulb[regular]}\ \textbf{Recuerda:}\ #1\end{tcolorbox}}
\newcommand{\res}[1]{\ifmmode\boxed{#1}\else\fbox{#1}\fi}

\def\fmneed{11}% que el títol d'una secció no quede sol al final del full
% =====================================================================
% Capçalera FentMates de les fitxes: un color per a cada curs.
% Els colors són els de la pàgina inicial (LEVELS a assets/js/materials.js);
% el logo ve de _eines/logo/pdf/ (python3 _eines/logo/genera.py).
%
% Ús, al final del preàmbul de la unitat:
%   % =====================================================================
% Capçalera FentMates de les fitxes: un color per a cada curs.
% Els colors són els de la pàgina inicial (LEVELS a assets/js/materials.js);
% el logo ve de _eines/logo/pdf/ (python3 _eines/logo/genera.py).
%
% Ús, al final del preàmbul de la unitat:
%   % =====================================================================
% Capçalera FentMates de les fitxes: un color per a cada curs.
% Els colors són els de la pàgina inicial (LEVELS a assets/js/materials.js);
% el logo ve de _eines/logo/pdf/ (python3 _eines/logo/genera.py).
%
% Ús, al final del preàmbul de la unitat:
%   \input{../comu/fentmates.tex}
%   \fmcurs{1eso}
% i, al començament del document:
%   \capcalera{Matemàtiques}{1r ESO}{Fitxa d'exercicis · Unitat 2}{Divisibilitat}
% Opcional abans de \input: \def\fmneed{14} (espai mínim abans de cada secció).
% =====================================================================
\makeatletter
\def\fm@logodir{../../logo/pdf}
\providecommand{\fmneed}{10}
\def\fm@cursnom{}

% clau del curs, color i color del text damunt del color
\newcommand{\fm@curs}[3]{\@namedef{fm@color@#1}{#2}\@namedef{fm@text@#1}{#3}}
\fm@curs{1eso}{E6B84A}{navy}
\fm@curs{2eso}{3B82F6}{white}
\fm@curs{3eso}{FF7A50}{white}
\fm@curs{4eso-a}{A78BFA}{white}
\fm@curs{4eso-b}{6D28D9}{white}
\fm@curs{1batx-mi}{9B1C1C}{white}
\fm@curs{1batx-ccss}{EF6B6B}{white}
\fm@curs{2batx-mii}{1E6B47}{white}
\fm@curs{2batx-ccss}{2BB3A3}{white}

\newcommand{\fmcurs}[1]{%
  \@ifundefined{fm@color@#1}{\PackageError{fentmates}{Curs desconegut: #1}{}}{}%
  \edef\fm@tmp{\noexpand\definecolor{curs}{HTML}{\@nameuse{fm@color@#1}}%
               \noexpand\colorlet{curstext}{\@nameuse{fm@text@#1}}}\fm@tmp
  \colorlet{cursfosc}{curs!72!black}%
  \def\fm@logo{\fm@logodir/vertical-#1.pdf}%
  \def\fm@marca{\fm@logodir/marca-#1.pdf}}

% ---------- capçalera de la primera pàgina ----------
% #1 matèria (petit), #2 curs, #3 tipus de document i unitat, #4 títol
\newcommand{\capcalera}[4]{%
  \gdef\fm@cursnom{#2}%
  \par\noindent
  \begin{tcolorbox}[colback=curs,colframe=curs,boxrule=0pt,arc=2.2mm,width=3.5cm,height=2.15cm,
      valign=center,left=3mm,right=2mm,top=1mm,bottom=1mm,nobeforeafter,box align=center]
    \color{curstext}{\footnotesize #1}\par\vspace{1pt}{\fontsize{19}{21}\selectfont\bfseries #2\par}%
  \end{tcolorbox}\hspace{5mm}%
  \begin{minipage}[c]{\dimexpr\linewidth-3.5cm-5mm-2.9cm\relax}
    {\color{cursfosc}\small\bfseries\MakeUppercase{#3}}\par\vspace{2pt}
    {\color{navy}\fontsize{20}{23}\selectfont\bfseries #4\par}
  \end{minipage}\hfill
  \begin{minipage}[c]{2.6cm}\raggedleft\includegraphics[width=2.6cm]{\fm@logo}\end{minipage}\par
  \vspace{2pt}}

% ---------- barres de secció del color del curs ----------
\renewcommand{\seccio}[2]{\par\needspace{\fmneed\baselineskip}\bigskip
  \begin{tcolorbox}[colback=curs,colframe=curs,coltext=curstext,arc=2mm,boxrule=0pt,left=3mm,right=3mm,top=1.2mm,bottom=1.2mm]
  {\large\bfseries #1\quad #2}
  \end{tcolorbox}\vspace{-2pt}}

% ---------- estrelles de dificultat del color del curs (a Batxillerat no se'n posen) ----------
\@ifundefined{estrelles}{}{%
  \renewcommand{\estrelles}[1]{{\color{curs}\ifcase#1\or\faStar\faStar[regular]\faStar[regular]\or\faStar\faStar\faStar[regular]\or\faStar\faStar\faStar\fi}}}

% ---------- capçalera i peu de les pàgines ----------
% primera pàgina: logos institucionals al peu
\fancypagestyle{primera}{\fancyhf{}\renewcommand{\headrulewidth}{0pt}
  \fancyfoot[L]{\raisebox{-0.32\height}{\includegraphics[height=0.78cm]{logo-gva.png}}\hspace{5mm}%
    \raisebox{-0.32\height}{\includegraphics[height=0.78cm]{logo-vermellar.png}}}
  \fancyfoot[R]{\scriptsize\color{muted}\href{https://creativecommons.org/licenses/by-nc-sa/4.0/deed.ca}{\ccbyncsa}\ \ CC BY-NC-SA 4.0 · fentmates.org\qquad\small\thepage}}
% la resta: el símbol a l'esquerra i el curs a la dreta, amb el títol al centre com abans
\fancyhead[L]{\raisebox{-0.3\height}{\includegraphics[height=0.62cm]{\fm@marca}}}
\fancyhead[R]{\small\bfseries\color{cursfosc}\fm@cursnom}
\renewcommand{\headrule}{\color{curs!45}\hrule height \headrulewidth}
\makeatother

%   \fmcurs{1eso}
% i, al començament del document:
%   \capcalera{Matemàtiques}{1r ESO}{Fitxa d'exercicis · Unitat 2}{Divisibilitat}
% Opcional abans de \input: \def\fmneed{14} (espai mínim abans de cada secció).
% =====================================================================
\makeatletter
\def\fm@logodir{../../logo/pdf}
\providecommand{\fmneed}{10}
\def\fm@cursnom{}

% clau del curs, color i color del text damunt del color
\newcommand{\fm@curs}[3]{\@namedef{fm@color@#1}{#2}\@namedef{fm@text@#1}{#3}}
\fm@curs{1eso}{E6B84A}{navy}
\fm@curs{2eso}{3B82F6}{white}
\fm@curs{3eso}{FF7A50}{white}
\fm@curs{4eso-a}{A78BFA}{white}
\fm@curs{4eso-b}{6D28D9}{white}
\fm@curs{1batx-mi}{9B1C1C}{white}
\fm@curs{1batx-ccss}{EF6B6B}{white}
\fm@curs{2batx-mii}{1E6B47}{white}
\fm@curs{2batx-ccss}{2BB3A3}{white}

\newcommand{\fmcurs}[1]{%
  \@ifundefined{fm@color@#1}{\PackageError{fentmates}{Curs desconegut: #1}{}}{}%
  \edef\fm@tmp{\noexpand\definecolor{curs}{HTML}{\@nameuse{fm@color@#1}}%
               \noexpand\colorlet{curstext}{\@nameuse{fm@text@#1}}}\fm@tmp
  \colorlet{cursfosc}{curs!72!black}%
  \def\fm@logo{\fm@logodir/vertical-#1.pdf}%
  \def\fm@marca{\fm@logodir/marca-#1.pdf}}

% ---------- capçalera de la primera pàgina ----------
% #1 matèria (petit), #2 curs, #3 tipus de document i unitat, #4 títol
\newcommand{\capcalera}[4]{%
  \gdef\fm@cursnom{#2}%
  \par\noindent
  \begin{tcolorbox}[colback=curs,colframe=curs,boxrule=0pt,arc=2.2mm,width=3.5cm,height=2.15cm,
      valign=center,left=3mm,right=2mm,top=1mm,bottom=1mm,nobeforeafter,box align=center]
    \color{curstext}{\footnotesize #1}\par\vspace{1pt}{\fontsize{19}{21}\selectfont\bfseries #2\par}%
  \end{tcolorbox}\hspace{5mm}%
  \begin{minipage}[c]{\dimexpr\linewidth-3.5cm-5mm-2.9cm\relax}
    {\color{cursfosc}\small\bfseries\MakeUppercase{#3}}\par\vspace{2pt}
    {\color{navy}\fontsize{20}{23}\selectfont\bfseries #4\par}
  \end{minipage}\hfill
  \begin{minipage}[c]{2.6cm}\raggedleft\includegraphics[width=2.6cm]{\fm@logo}\end{minipage}\par
  \vspace{2pt}}

% ---------- barres de secció del color del curs ----------
\renewcommand{\seccio}[2]{\par\needspace{\fmneed\baselineskip}\bigskip
  \begin{tcolorbox}[colback=curs,colframe=curs,coltext=curstext,arc=2mm,boxrule=0pt,left=3mm,right=3mm,top=1.2mm,bottom=1.2mm]
  {\large\bfseries #1\quad #2}
  \end{tcolorbox}\vspace{-2pt}}

% ---------- estrelles de dificultat del color del curs (a Batxillerat no se'n posen) ----------
\@ifundefined{estrelles}{}{%
  \renewcommand{\estrelles}[1]{{\color{curs}\ifcase#1\or\faStar\faStar[regular]\faStar[regular]\or\faStar\faStar\faStar[regular]\or\faStar\faStar\faStar\fi}}}

% ---------- capçalera i peu de les pàgines ----------
% primera pàgina: logos institucionals al peu
\fancypagestyle{primera}{\fancyhf{}\renewcommand{\headrulewidth}{0pt}
  \fancyfoot[L]{\raisebox{-0.32\height}{\includegraphics[height=0.78cm]{logo-gva.png}}\hspace{5mm}%
    \raisebox{-0.32\height}{\includegraphics[height=0.78cm]{logo-vermellar.png}}}
  \fancyfoot[R]{\scriptsize\color{muted}\href{https://creativecommons.org/licenses/by-nc-sa/4.0/deed.ca}{\ccbyncsa}\ \ CC BY-NC-SA 4.0 · fentmates.org\qquad\small\thepage}}
% la resta: el símbol a l'esquerra i el curs a la dreta, amb el títol al centre com abans
\fancyhead[L]{\raisebox{-0.3\height}{\includegraphics[height=0.62cm]{\fm@marca}}}
\fancyhead[R]{\small\bfseries\color{cursfosc}\fm@cursnom}
\renewcommand{\headrule}{\color{curs!45}\hrule height \headrulewidth}
\makeatother

%   \fmcurs{1eso}
% i, al començament del document:
%   \capcalera{Matemàtiques}{1r ESO}{Fitxa d'exercicis · Unitat 2}{Divisibilitat}
% Opcional abans de \input: \def\fmneed{14} (espai mínim abans de cada secció).
% =====================================================================
\makeatletter
\def\fm@logodir{../../logo/pdf}
\providecommand{\fmneed}{10}
\def\fm@cursnom{}

% clau del curs, color i color del text damunt del color
\newcommand{\fm@curs}[3]{\@namedef{fm@color@#1}{#2}\@namedef{fm@text@#1}{#3}}
\fm@curs{1eso}{E6B84A}{navy}
\fm@curs{2eso}{3B82F6}{white}
\fm@curs{3eso}{FF7A50}{white}
\fm@curs{4eso-a}{A78BFA}{white}
\fm@curs{4eso-b}{6D28D9}{white}
\fm@curs{1batx-mi}{9B1C1C}{white}
\fm@curs{1batx-ccss}{EF6B6B}{white}
\fm@curs{2batx-mii}{1E6B47}{white}
\fm@curs{2batx-ccss}{2BB3A3}{white}

\newcommand{\fmcurs}[1]{%
  \@ifundefined{fm@color@#1}{\PackageError{fentmates}{Curs desconegut: #1}{}}{}%
  \edef\fm@tmp{\noexpand\definecolor{curs}{HTML}{\@nameuse{fm@color@#1}}%
               \noexpand\colorlet{curstext}{\@nameuse{fm@text@#1}}}\fm@tmp
  \colorlet{cursfosc}{curs!72!black}%
  \def\fm@logo{\fm@logodir/vertical-#1.pdf}%
  \def\fm@marca{\fm@logodir/marca-#1.pdf}}

% ---------- capçalera de la primera pàgina ----------
% #1 matèria (petit), #2 curs, #3 tipus de document i unitat, #4 títol
\newcommand{\capcalera}[4]{%
  \gdef\fm@cursnom{#2}%
  \par\noindent
  \begin{tcolorbox}[colback=curs,colframe=curs,boxrule=0pt,arc=2.2mm,width=3.5cm,height=2.15cm,
      valign=center,left=3mm,right=2mm,top=1mm,bottom=1mm,nobeforeafter,box align=center]
    \color{curstext}{\footnotesize #1}\par\vspace{1pt}{\fontsize{19}{21}\selectfont\bfseries #2\par}%
  \end{tcolorbox}\hspace{5mm}%
  \begin{minipage}[c]{\dimexpr\linewidth-3.5cm-5mm-2.9cm\relax}
    {\color{cursfosc}\small\bfseries\MakeUppercase{#3}}\par\vspace{2pt}
    {\color{navy}\fontsize{20}{23}\selectfont\bfseries #4\par}
  \end{minipage}\hfill
  \begin{minipage}[c]{2.6cm}\raggedleft\includegraphics[width=2.6cm]{\fm@logo}\end{minipage}\par
  \vspace{2pt}}

% ---------- barres de secció del color del curs ----------
\renewcommand{\seccio}[2]{\par\needspace{\fmneed\baselineskip}\bigskip
  \begin{tcolorbox}[colback=curs,colframe=curs,coltext=curstext,arc=2mm,boxrule=0pt,left=3mm,right=3mm,top=1.2mm,bottom=1.2mm]
  {\large\bfseries #1\quad #2}
  \end{tcolorbox}\vspace{-2pt}}

% ---------- estrelles de dificultat del color del curs (a Batxillerat no se'n posen) ----------
\@ifundefined{estrelles}{}{%
  \renewcommand{\estrelles}[1]{{\color{curs}\ifcase#1\or\faStar\faStar[regular]\faStar[regular]\or\faStar\faStar\faStar[regular]\or\faStar\faStar\faStar\fi}}}

% ---------- capçalera i peu de les pàgines ----------
% primera pàgina: logos institucionals al peu
\fancypagestyle{primera}{\fancyhf{}\renewcommand{\headrulewidth}{0pt}
  \fancyfoot[L]{\raisebox{-0.32\height}{\includegraphics[height=0.78cm]{logo-gva.png}}\hspace{5mm}%
    \raisebox{-0.32\height}{\includegraphics[height=0.78cm]{logo-vermellar.png}}}
  \fancyfoot[R]{\scriptsize\color{muted}\href{https://creativecommons.org/licenses/by-nc-sa/4.0/deed.ca}{\ccbyncsa}\ \ CC BY-NC-SA 4.0 · fentmates.org\qquad\small\thepage}}
% la resta: el símbol a l'esquerra i el curs a la dreta, amb el títol al centre com abans
\fancyhead[L]{\raisebox{-0.3\height}{\includegraphics[height=0.62cm]{\fm@marca}}}
\fancyhead[R]{\small\bfseries\color{cursfosc}\fm@cursnom}
\renewcommand{\headrule}{\color{curs!45}\hrule height \headrulewidth}
\makeatother

\fmcurs{4eso-a}

\begin{document}
\thispagestyle{primera}

\capcalera{Matemáticas A}{4\textsuperscript{o} ESO}{Soluciones de la ficha · Unidad 1}{\raggedright Números naturales, enteros y fraccionarios}
\vspace{8pt}

\begin{tcolorbox}[colback=soft,colframe=soft,arc=2mm,left=3mm,right=3mm,top=1.5mm,bottom=1.5mm]
\small Intenta resolver cada ejercicio antes de mirar su solución. Si no te sale el mismo resultado, repasa los pasos o consulta los apuntes interactivos de la unidad.
\end{tcolorbox}

% =====================================================
\seccio{1}{Números naturales. Operaciones combinadas}

\ex{1} \begin{tasks}(3)
  \task $3^{5}=\res{243}$
  \task $10^{6}=\res{1000000}$
  \task $2^{3}\cdot 5^{2}=\res{200}$
  \task $\sqrt{144}=\res{12}$
  \task $\sqrt{225}=\res{15}$
  \task $17^{0}=\res{1}$
\end{tasks}

\ex{1} \begin{tasks}(2)
  \task $1250=37\cdot\res{33}+\res{29}$
  \task $5000=123\cdot\res{40}+\res{80}$
\end{tasks}

\ex{2} \begin{tasks}(1)
  \task $25-3\cdot \left[12-2\cdot \left(7-4\right)\right]=25-3\cdot \left[12-2\cdot 3\right]=25-3\cdot 6=25-18=\res{7}$
  \task $6+48:\left[20-4\cdot \left(9-6\right)\right]=6+48:\left[20-4\cdot 3\right]=6+48:8=6+6=\res{12}$
  \task $60:\left[3\cdot \left(9-5\right)+2\cdot \left(7-3\right)\right]=60:\left[3\cdot 4+2\cdot 4\right]=60:\left[12+8\right]=60:20=\res{3}$
  \task $2\cdot \left(40-3\cdot \left[5+2\cdot \left(9-6\right)\right]\right)+7=2\cdot \left(40-3\cdot \left[5+2\cdot 3\right]\right)+7=2\cdot \left(40-3\cdot 11\right)+7=2\cdot 7+7=\res{21}$
  \task $5\cdot \left[18-\left(4+2\cdot 3\right)\right]-4\cdot \left[10-\left(9-2\cdot 3\right)\right]=5\cdot \left[18-10\right]-4\cdot \left[10-3\right]=5\cdot 8-4\cdot 7=40-28=\res{12}$
  \task $2^{3}\cdot \left[15-\left(\sqrt{49}+2\right)\right]+36:3^{2}=8\cdot \left[15-9\right]+36:9=8\cdot 6+4=48+4=\res{52}$
\end{tasks}

\ex{2} $2\cdot 100-\left(3\cdot 18+2\cdot 35-15\right)=200-\left(54+70-15\right)=200-109=\res{91}$\,€

\ex{3} Ha multiplicado antes de dividir: las multiplicaciones y divisiones se hacen de izquierda a derecha. $36:6\cdot 2+4^{2}=6\cdot 2+16=12+16=\res{28}$

\ex{3} Lado: $\sqrt{1225}=35$\,m. Una vuelta: $4\cdot35-5=135$\,m; tres vueltas: $3\cdot135=405$\,m. $405:50=8$ y sobran 5\,m, por tanto hacen falta $\res{9\text{ rollos}}$. Coste: $9\cdot18=\res{162\text{ €}}$.


% =====================================================
\seccio{2}{Números enteros. Valor absoluto}

\ex{1} $\res{-21<-15<-8<-3<0<5<8<12}$

\ex{1} \begin{tasks}(3)
  \task $\left\lvert -15\right\rvert +\left\lvert 7\right\rvert =\res{22}$
  \task $\left\lvert -4-9\right\rvert =\res{13}$
  \task $\left\lvert -4\right\rvert -\left\lvert 9\right\rvert =\res{-5}$
  \task $-\left\lvert -12\right\rvert =\res{-12}$
  \task $\left\lvert 3-11\right\rvert \cdot \left\lvert -2\right\rvert =\res{16}$
  \task $\left\lvert -20\right\rvert :\left\lvert -5\right\rvert -\left\lvert -3\right\rvert =\res{1}$
\end{tasks}

\ex{2} \begin{tasks}(1)
  \task $x=\res{-9}$ y $x=\res{9}$
  \task $\res{-3,\ {-2},\ {-1},\ 0,\ 1,\ 2,\ 3}$
  \task $\res{-4,\ {-3},\ {-2},\ {-1},\ 0,\ 1}$
  \task $x$ está a 5 unidades del 2: \ $x=\res{-3}$ y $x=\res{7}$
\end{tasks}

\ex{3} \begin{tasks}(1)
  \task $\left\lvert 4-9\right\rvert -\left\lvert 2-3\cdot 5\right\rvert +\left\lvert -6\right\rvert \cdot \left(-2\right)=\left\lvert -5\right\rvert -\left\lvert -13\right\rvert +6\cdot \left(-2\right)=5-13-12=\res{-20}$
  \task $\left(3-8\right)\cdot 2+\left\lvert 6-4\cdot 3\right\rvert \cdot \left\lvert -1-2\right\rvert =\left(-5\right)\cdot 2+\left\lvert -6\right\rvert \cdot \left\lvert -3\right\rvert =-10+6\cdot 3=-10+18=\res{8}$
  \task $7-\left\lvert -3-\left\lvert 2-9\right\rvert \right\rvert =7-\left\lvert -3-7\right\rvert =7-\left\lvert -10\right\rvert =7-10=\res{-3}$
  \task $\left\lvert \left(-2\right)\cdot \left(-5\right)-\left\lvert 3-15\right\rvert \right\rvert -\left\lvert -4\right\rvert \cdot \left(1-3\right)=\left\lvert 10-12\right\rvert -4\cdot \left(-2\right)=2+8=\res{10}$
\end{tasks}

\ex{2} \begin{tasks}(2)
  \task $-27+19=\res{-8}$\,°C; \ $-8-14=\res{-22}$\,°C
  \task $-8-\left(-27\right)=-8+27=\res{19}$\,°C
\end{tasks}

\ex{2} $-212-\left(-287\right)=-212+287=\res{75}$ años.


% =====================================================
\seccio{3}{Operaciones combinadas con enteros}

\ex{1} \begin{tasks}(3)
  \task $-12+7-\left(-5\right)=\res{0}$
  \task $\left(-8\right)\cdot \left(-6\right):\left(-4\right)=\res{-12}$
  \task $\left(-72\right):\left(-9\right)\cdot \left(-2\right)=\res{-16}$
  \task $-15-\left(-3\right)\cdot 4=\res{-3}$
  \task $\left(-2\right)\cdot \left(-3\right)\cdot \left(-5\right)=\res{-30}$
  \task $-36:\left(-4\right)-9=\res{0}$
\end{tasks}

\ex{2} \begin{tasks}(1)
  \task $12-4\cdot \left[6-3\cdot \left(5-8\right)\right]=12-4\cdot \left[6-3\cdot \left(-3\right)\right]=12-4\cdot \left[6+9\right]=12-4\cdot 15=12-60=\res{-48}$
  \task $-3+24:\left[-2-2\cdot \left(1-6\right)\right]=-3+24:\left[-2-2\cdot \left(-5\right)\right]=-3+24:\left[-2+10\right]=-3+24:8=-3+3=\res{0}$
  \task $5\cdot \left(-4+9\right)-\left[7-\left(-3\right)\cdot \left(2-6\right)\right]=5\cdot 5-\left[7-\left(-3\right)\cdot \left(-4\right)\right]=25-\left[7-12\right]=25-\left(-5\right)=\res{30}$
  \task $\left(-2\right)\cdot \left[8-\left(3-5\right)\cdot \left(-4\right)\right]-\left(-6\right):\left(-2\right)=\left(-2\right)\cdot \left[8-\left(-2\right)\cdot \left(-4\right)\right]-3=\left(-2\right)\cdot \left[8-8\right]-3=\left(-2\right)\cdot 0-3=\res{-3}$
  \task $40:\left[-12+4\cdot \left(7-9\right)\cdot \left(-1\right)\right]=40:\left[-12+4\cdot \left(-2\right)\cdot \left(-1\right)\right]=40:\left[-12+8\right]=40:\left(-4\right)=\res{-10}$
  \task $-\left[6-\left(2-11\right)\right]+3\cdot \left[-4+\left(5-1\right)\cdot 2\right]=-\left[6-\left(-9\right)\right]+3\cdot \left[-4+4\cdot 2\right]=-15+3\cdot 4=-15+12=\res{-3}$
\end{tasks}

\ex{2} \begin{tasks}(1)
  \task $\left(-3\right)^{3}+\left(-2\right)^{4}-\left(-1\right)^{9}=-27+16-\left(-1\right)=-27+16+1=\res{-10}$
  \task $-4^{2}+\left(-4\right)^{2}-\sqrt{36}\cdot \left(-2\right)^{3}=-16+16-6\cdot \left(-8\right)=0+48=\res{48}$
  \task $\left[\left(-5\right)^{2}-\sqrt{100-36}\cdot 3\right]:\left(-1\right)^{7}=\left[25-\sqrt{64}\cdot 3\right]:\left(-1\right)=\left[25-24\right]:\left(-1\right)=1:\left(-1\right)=\res{-1}$
  \task $2\cdot \left(-3\right)^{2}-\left[\left(-2\right)^{3}+\sqrt{81}\right]^{2}=2\cdot 9-\left[-8+9\right]^{2}=18-1^{2}=18-1=\res{17}$
  \task $\left(-2\right)^{5}:\left(-2\right)^{3}-3\cdot \left(-1\right)^{20}+\sqrt{4^{2}+3^{2}}=\left(-32\right):\left(-8\right)-3\cdot 1+\sqrt{25}=4-3+5=\res{6}$
\end{tasks}

\ex{2} $420-3\cdot 85-2\cdot 120+2\cdot 150=420-255-240+300=\res{225}$\,€. Sí: después de los recibos el saldo era $420-495=\res{-75}$\,€.

\ex{3} $3-5+2\cdot 5-9=3-5+10-9=\res{-1}$ (primer sótano). \quad Recorrido: $\left(5+10+9\right)\cdot 3=24\cdot 3=\res{72}$\,m

\ex{3} Media: $\left(-6+3-2-8+1-4+2\right):7=\left(-14\right):7=\res{-2}$\,°C. \quad Diferencia: $3-\left(-8\right)=\res{11}$\,°C.

\ex{3} Han de tener signos distintos y el negativo debe ser el de valor absoluto mayor: $\res{4}$ y $\res{-6}$ \ ($4\cdot(-6)=-24$, \ $4+(-6)=-2$).


% =====================================================
\seccio{4}{Fracciones. Fracciones equivalentes}

\ex{1} \begin{tasks}[after-item-skip=6pt](2)
  \task $\frac{84}{140}=\frac{84:28}{140:28}=\res{\frac{3}{5}}$
  \task $\frac{126}{315}=\frac{126:63}{315:63}=\res{\frac{2}{5}}$
  \task $\frac{210}{462}=\frac{210:42}{462:42}=\res{\frac{5}{11}}$
  \task $\frac{360}{840}=\frac{360:120}{840:120}=\res{\frac{3}{7}}$
  \task $\frac{156}{234}=\frac{156:78}{234:78}=\res{\frac{2}{3}}$
  \task $\frac{275}{1100}=\frac{275:275}{1100:275}=\res{\frac{1}{4}}$
\end{tasks}

\ex{1} \begin{tasks}(1)
  \task $\frac{2}{3}\cdot\frac{3}{5}\cdot450=\frac{2}{5}\cdot450=\res{180\text{ €}}$
  \task $\frac{3}{4}\cdot\frac{1}{2}\cdot96=\frac{3}{8}\cdot96=\res{36}$
  \task $\frac{5}{6}\cdot\frac{1}{3}\cdot720=\frac{5}{18}\cdot720=\res{200\text{ L}}$
\end{tasks}

\ex{2} \begin{tasks}(1)
  \task $21\cdot55=1155=35\cdot33$: \res{sí}
  \task $18\cdot32=576\neq24\cdot27=648$: \res{no}
  \task $(-14)\cdot(-21)=294=49\cdot6$: \res{sí}
\end{tasks}

\ex{2} \begin{tasks}(3)
  \task $42x=630\Rightarrow x=\res{15}$
  \task $-4x=84\Rightarrow x=\res{-21}$
  \task $25x=600\Rightarrow x=\res{24}$
\end{tasks}

\ex{2} $\frac{16}{24}=\res{\frac{2}{3}}$; \quad $\frac{45}{1440}=\res{\frac{1}{32}}$; \quad $\frac{160}{1440}=\res{\frac{1}{9}}$

\ex{2} Con denominador común 36: $\frac{30}{36},\ -\frac{27}{36},\ \frac{28}{36},\ -\frac{24}{36},\ \frac{26}{36}$. \quad $\res{-\frac{3}{4}<-\frac{2}{3}<0<\frac{13}{18}<\frac{7}{9}<\frac{5}{6}}$

\ex{3} Amplificamos para que quede espacio entre los numeradores: $\frac{6}{14}<\res{\frac{7}{14}}<\frac{8}{14}$; \ $\frac{8}{20}<\res{\frac{9}{20}}<\frac{10}{20}$ (por ejemplo).


% =====================================================
\seccio{5}{Operaciones con fracciones}

\ex{1} \begin{tasks}[after-item-skip=6pt](2)
  \task $\frac{5}{12}+\frac{7}{18}=\frac{15}{36}+\frac{14}{36}=\res{\frac{29}{36}}$
  \task $\frac{11}{15}-\frac{7}{10}=\frac{22}{30}-\frac{21}{30}=\res{\frac{1}{30}}$
  \task $3-\frac{11}{4}+\frac{5}{6}=\frac{36}{12}-\frac{33}{12}+\frac{10}{12}=\res{\frac{13}{12}}$
  \task $-\frac{7}{20}+\frac{5}{8}=-\frac{14}{40}+\frac{25}{40}=\res{\frac{11}{40}}$
\end{tasks}

\ex{1} \begin{tasks}[after-item-skip=6pt](2)
  \task $\frac{14}{15}\cdot \frac{25}{21}=\frac{350}{315}=\res{\frac{10}{9}}$
  \task $-\frac{9}{16}\cdot \frac{8}{27}=-\frac{72}{432}=\res{-\frac{1}{6}}$
  \task $\frac{12}{35}:\frac{18}{25}=\frac{300}{630}=\res{\frac{10}{21}}$
  \task $\left(-\frac{5}{6}\right):\left(-\frac{10}{9}\right)=\frac{45}{60}=\res{\frac{3}{4}}$
\end{tasks}

\ex{2} \begin{tasks}[after-item-skip=6pt](2)
  \task $\frac{2}{3}\cdot \frac{9}{10}=\res{\frac{3}{5}}$
  \task $\frac{1}{2}\cdot \frac{4}{7}=\res{\frac{2}{7}}$
  \task $\frac{3}{5}\cdot \frac{5}{12}=\res{\frac{1}{4}}$
  \task $\frac{7}{8}\cdot \frac{4}{21}=\res{\frac{1}{6}}$
\end{tasks}

\ex{2} \begin{tasks}[after-item-skip=6pt](2)
  \task $-\frac{5}{6}-\left(-\frac{3}{4}\right)+-\frac{1}{3}=-\frac{10}{12}+\frac{9}{12}-\frac{4}{12}=\res{-\frac{5}{12}}$
  \task $\left(-\frac{2}{3}\right)\cdot \frac{9}{4}\cdot \left(-\frac{8}{15}\right)=\frac{144}{180}=\res{\frac{4}{5}}$
  \task $\left(-\frac{4}{9}\right):\left(-\frac{2}{3}\right)\cdot \left(-6\right)=\frac{2}{3}\cdot \left(-6\right)=\res{-4}$
  \task $-2-\left(-\frac{5}{4}\right)+\frac{7}{6}=-\frac{24}{12}+\frac{15}{12}+\frac{14}{12}=\res{\frac{5}{12}}$
\end{tasks}

\ex{2} \begin{tasks}(2)
  \task Op. $\res{-\frac{3}{8}}$, \ inv. $\res{\frac{8}{3}}$
  \task Op. $\res{\frac{5}{2}}$, \ inv. $\res{-\frac{2}{5}}$
  \task Op. $\res{6}$, \ inv. $\res{-\frac{1}{6}}$
  \task $1+\frac{1}{3}=\frac{4}{3}$: \ op. $\res{-\frac{4}{3}}$, \ inv. $\res{\frac{3}{4}}$
\end{tasks}


% =====================================================
\seccio{6}{Operaciones combinadas con fracciones}

\ex{1} \begin{tasks}[after-item-skip=6pt](2)
  \task $\frac{2}{3}+\frac{1}{2}\cdot \frac{4}{9}=\frac{2}{3}+\frac{2}{9}=\frac{6}{9}+\frac{2}{9}=\res{\frac{8}{9}}$
  \task $\left(\frac{2}{3}+\frac{1}{2}\right)\cdot \frac{4}{9}=\frac{7}{6}\cdot \frac{4}{9}=\frac{28}{54}=\res{\frac{14}{27}}$
  \task $\frac{3}{4}-\frac{1}{4}:\frac{3}{8}=\frac{3}{4}-\frac{2}{3}=\frac{9}{12}-\frac{8}{12}=\res{\frac{1}{12}}$
  \task $\left(\frac{3}{4}-\frac{1}{4}\right):\frac{3}{8}=\frac{1}{2}:\frac{3}{8}=\frac{8}{6}=\res{\frac{4}{3}}$
\end{tasks}

\ex{2} \begin{tasks}[after-item-skip=6pt](1)
  \task $\frac{3}{4}-\left[\frac{1}{2}-\left(\frac{2}{5}-\frac{1}{10}\right)\right]=\frac{3}{4}-\left[\frac{1}{2}-\frac{3}{10}\right]=\frac{3}{4}-\frac{1}{5}=\frac{15}{20}-\frac{4}{20}=\res{\frac{11}{20}}$
  \task $2-\left[\frac{1}{3}+\left(\frac{5}{6}-1\right)\right]+\frac{1}{4}=2-\left[\frac{1}{3}+\left(-\frac{1}{6}\right)\right]+\frac{1}{4}=2-\frac{1}{6}+\frac{1}{4}=\frac{24}{12}-\frac{2}{12}+\frac{3}{12}=\res{\frac{25}{12}}$
  \task $1-\left(\frac{2}{3}+\frac{1}{4}-\frac{1}{2}\right):\frac{5}{6}=1-\frac{5}{12}:\frac{5}{6}=1-\frac{1}{2}=\res{\frac{1}{2}}$
  \task $\frac{4}{5}\cdot \left[\frac{1}{2}-\left(\frac{3}{4}-\frac{1}{8}\right)\right]=\frac{4}{5}\cdot \left[\frac{1}{2}-\frac{5}{8}\right]=\frac{4}{5}\cdot \left(-\frac{1}{8}\right)=\res{-\frac{1}{10}}$
  \task $\left[\frac{2}{3}\cdot \left(\frac{1}{4}+\frac{1}{2}\right)\right]:\left[\left(-3\right)\cdot \left(\frac{1}{6}-\frac{1}{2}\right)\right]=\left[\frac{2}{3}\cdot \frac{3}{4}\right]:\left[\left(-3\right)\cdot \left(-\frac{1}{3}\right)\right]=\frac{1}{2}:1=\res{\frac{1}{2}}$
  \task $\left(\frac{5}{6}-\frac{2}{3}-\frac{1}{12}\right):\left(1-\frac{3}{4}\right)=\frac{1}{12}:\frac{1}{4}=\res{\frac{1}{3}}$
\end{tasks}

\ex{2} \begin{tasks}[after-item-skip=6pt](1)
  \task $\left(\frac{3}{2}\right)^{2}-\frac{5}{4}\cdot \left(1-\frac{1}{5}\right)=\frac{9}{4}-\frac{5}{4}\cdot \frac{4}{5}=\frac{9}{4}-1=\res{\frac{5}{4}}$
  \task $\left(-\frac{2}{3}\right)^{3}:\frac{4}{9}+\frac{1}{2}=-\frac{8}{27}:\frac{4}{9}+\frac{1}{2}=-\frac{2}{3}+\frac{1}{2}=\res{-\frac{1}{6}}$
  \task $\left[1-\left(\frac{1}{2}\right)^{2}\right]^{2}\cdot \frac{16}{3}=\left[1-\frac{1}{4}\right]^{2}\cdot \frac{16}{3}=\left(\frac{3}{4}\right)^{2}\cdot \frac{16}{3}=\frac{9}{16}\cdot \frac{16}{3}=\res{3}$
\end{tasks}

\ex{3} \begin{tasks}[after-item-skip=10pt](2)
  \task $\dfrac{\frac{2}{3}-\frac{1}{4}}{\frac{5}{6}+\frac{1}{3}}=\dfrac{\frac{5}{12}}{\frac{7}{6}}=\res{\frac{5}{14}}$
  \task $\dfrac{1-\frac{3}{5}\cdot \frac{5}{6}}{\left(\frac{1}{2}\right)^{2}+\frac{1}{3}}=\dfrac{\frac{1}{2}}{\frac{7}{12}}=\res{\frac{6}{7}}$
  \task $\dfrac{\left(-2\right)\cdot \left(\frac{1}{4}-\frac{2}{3}\right)}{3\cdot \left(\frac{5}{6}-\frac{1}{2}\right)}=\dfrac{\left(-2\right)\cdot \left(-\frac{5}{12}\right)}{3\cdot \frac{1}{3}}=\dfrac{\frac{5}{6}}{1}=\res{\frac{5}{6}}$
  \task $\dfrac{\frac{3}{4}:\frac{9}{8}-\frac{1}{3}}{1-\frac{4}{9}}=\dfrac{\frac{2}{3}-\frac{1}{3}}{\frac{5}{9}}=\dfrac{\frac{1}{3}}{\frac{5}{9}}=\res{\frac{3}{5}}$
\end{tasks}

\ex{3} \begin{tasks}[after-item-skip=6pt](1)
  \task $\left[\frac{3}{4}-\left(\frac{1}{2}\right)^{2}\right]^{2}:\left(-\frac{1}{2}\right)^{3}+\frac{5}{3}\cdot \frac{6}{5}=\left[\frac{3}{4}-\frac{1}{4}\right]^{2}:\left(-\frac{1}{8}\right)+2=\left(\frac{1}{2}\right)^{2}:\left(-\frac{1}{8}\right)+2=\frac{1}{4}:\left(-\frac{1}{8}\right)+2=-2+2=\res{0}$
  \task $\frac{1}{2}-\frac{2}{3}\cdot \left[\frac{3}{4}-\left(\frac{1}{2}-2\right)^{2}:3\right]=\frac{1}{2}-\frac{2}{3}\cdot \left[\frac{3}{4}-\left(-\frac{3}{2}\right)^{2}:3\right]=\frac{1}{2}-\frac{2}{3}\cdot \left[\frac{3}{4}-\frac{9}{4}:3\right]=\frac{1}{2}-\frac{2}{3}\cdot \left[\frac{3}{4}-\frac{3}{4}\right]=\frac{1}{2}-\frac{2}{3}\cdot 0=\res{\frac{1}{2}}$
\end{tasks}


% =====================================================
\seccio{7}{Problemas con fracciones}

\ex{1} $352:\frac{4}{9}=\frac{352\cdot9}{4}=\res{792}$ alumnos.

\ex{2} Entrada: $\frac{1}{4}\cdot480=120$\,€; resto: 360\,€. Cuota: $360:9=\res{40\text{ €}}$, \ que es $\frac{3}{4}:9=\res{\frac{1}{12}}$ del precio.

\ex{2} 1.º: $\frac{2}{5}\cdot1800=\res{720}$\,€; \ quedan 1080\,€. \ 2.º: $\frac{3}{4}\cdot1080=\res{810}$\,€. \ 3.º: $1080-810=\res{270}$\,€, \ $\frac{270}{1800}=\res{\frac{3}{20}}$.

\ex{2} Tomates: $\frac{3}{8}\cdot2400=\res{900}$\,m²; \ resto 1500\,m². \ Patatas: $\frac{2}{5}\cdot1500=\res{600}$\,m². \ Pimientos: $\res{900}$\,m², \ $\frac{900}{2400}=\res{\frac{3}{8}}$.

\ex{3} $\frac{1}{6}+\frac{1}{4}=\frac{2}{12}+\frac{3}{12}=\res{\frac{5}{12}}$. \ Tiempo: $1:\frac{5}{12}=\frac{12}{5}=2{,}4$\,h $=\res{2\text{ h }24\text{ min}}$.

\ex{3} Después del alquiler le quedan $\frac{5}{7}$; en comida gasta $\frac{1}{3}\cdot\frac{5}{7}=\frac{5}{21}$. Le quedan $\frac{5}{7}-\frac{5}{21}=\frac{10}{21}$, que son 600\,€: \ $600:\frac{10}{21}=\res{1260\text{ €}}$.

\ex{3} Queda $\frac{1}{4}$ de $\frac{3}{5}$, es decir, $\frac{1}{4}\cdot\frac{3}{5}=\frac{3}{20}$ del viaje: \ $90:\frac{3}{20}=\res{600\text{ km}}$.

\ex{3} Los $\frac{3}{5}$ de la miel pesan $950-470=480$\,g; toda la miel, $480:\frac{3}{5}=800$\,g. Bote: $950-800=\res{150\text{ g}}$.

\ex{3} Queda la fracción $\frac{2}{3}\cdot\frac{3}{4}\cdot\frac{3}{5}=\frac{18}{60}=\frac{3}{10}$ del depósito: \ $360:\frac{3}{10}=\res{1200\text{ L}}$.


% =====================================================
\seccio{8}{Potencias con base entera y racional}

\ex{1} \begin{tasks}(3)
  \task $\left(-4\right)^{3}=\res{-64}$
  \task $-4^{2}=\res{-16}$
  \task $\left(-1\right)^{23}=\res{-1}$
  \task $-\left(-2\right)^{4}=\res{-16}$
  \task $\left(-10\right)^{5}=\res{-100000}$
  \task $\left(-9\right)^{0}=\res{1}$
\end{tasks}

\ex{1} \begin{tasks}[after-item-skip=6pt](3)
  \task $\left(\frac{3}{4}\right)^{3}=\res{\frac{27}{64}}$
  \task $\left(-\frac{2}{3}\right)^{4}=\res{\frac{16}{81}}$
  \task $\left(-\frac{1}{5}\right)^{3}=\res{-\frac{1}{125}}$
  \task $-\left(\frac{3}{2}\right)^{2}=\res{-\frac{9}{4}}$
  \task $\left(-\frac{7}{8}\right)^{0}=\res{1}$
  \task $\left(-\frac{1}{10}\right)^{5}=\res{-\frac{1}{100000}}$
\end{tasks}

\ex{2} \begin{tasks}(5)
  \task $\res{-}$
  \task $\res{+}$
  \task $\res{-}$
  \task $\res{-}$
  \task $\res{-}$
\end{tasks}

\ex{2} \begin{tasks}[after-item-skip=6pt](2)
  \task $\left(-3\right)^{3}-\left(-3\right)^{2}-3^{2}=-27-9-9=\res{-45}$
  \task $\left(-1\right)^{15}-\left(-1\right)^{16}+\left(-2\right)^{3}=-1-1+\left(-8\right)=\res{-10}$
  \task $\left(\frac{2}{3}\right)^{2}-\left(\frac{1}{3}\right)^{2}=\frac{4}{9}-\frac{1}{9}=\frac{3}{9}=\res{\frac{1}{3}}$
  \task $\left(-\frac{1}{2}\right)^{3}\cdot 8-\left(\frac{3}{4}\right)^{2}:\frac{9}{8}=-\frac{1}{8}\cdot 8-\frac{9}{16}:\frac{9}{8}=-1-\frac{1}{2}=\res{-\frac{3}{2}}$
\end{tasks}

\ex{3} a) $16\cdot \left(\frac{3}{4}\right)^{4}=16\cdot \frac{81}{256}=\res{\frac{81}{16}}$\,m $\approx5{,}06$\,m. \quad b) Alturas de los botes: 12\,m i 9\,m; \ $16+2\cdot 12+2\cdot 9=16+24+18=\res{58}$\,m.


% =====================================================
\seccio{9}{Propiedades de las potencias}

\ex{1} \begin{tasks}(3)
  \task $5^{3}\cdot 5^{4}=\res{5^{7}}$
  \task $2^{9}:2^{4}=\res{2^{5}}$
  \task $\left(3^{2}\right)^{5}=\res{3^{10}}$
  \task $4^{3}\cdot 5^{3}=\res{20^{3}}$
  \task $18^{6}:6^{6}=\res{3^{6}}$
  \task $\left(-7\right)^{4}\cdot \left(-7\right)^{3}=\res{(-7)^{7}}$
\end{tasks}

\ex{1} \begin{tasks}[after-item-skip=6pt](2)
  \task $\left(\frac{2}{5}\right)^3\cdot\left(\frac{2}{5}\right)^4=\res{\left(\frac{2}{5}\right)^{7}}$
  \task $\left(\frac{3}{4}\right)^8:\left(\frac{3}{4}\right)^5=\res{\left(\frac{3}{4}\right)^{3}}$
  \task $\left[\left(-\frac{1}{2}\right)^3\right]^2=\res{\left(-\frac{1}{2}\right)^{6}}$
  \task $\left(\frac{2}{3}\right)^4\cdot\left(\frac{9}{4}\right)^4=\res{\left(\frac{3}{2}\right)^{4}}$
\end{tasks}

\ex{2} \begin{tasks}[after-item-skip=6pt](1)
  \task $\left(-2\right)^{7}:\left(-2\right)^{4}=\left(-2\right)^{3}=\res{-8}$
  \task $\left[\left(-1\right)^{3}\right]^{5}\cdot \left(-1\right)^{2}=\left(-1\right)^{15}\cdot \left(-1\right)^{2}=\left(-1\right)^{17}=\res{-1}$
  \task $2^{5}\cdot 2^{3}:2^{6}=2^{8}:2^{6}=2^{2}=\res{4}$
  \task $\left(\frac{2}{3}\right)^{5}\cdot \left(\frac{2}{3}\right)^{2}:\left(\frac{2}{3}\right)^{4}=\left(\frac{2}{3}\right)^{7}:\left(\frac{2}{3}\right)^{4}=\left(\frac{2}{3}\right)^{3}=\res{\frac{8}{27}}$
  \task $\left(-\frac{1}{3}\right)^{6}:\left(-\frac{1}{3}\right)^{4}=\left(-\frac{1}{3}\right)^{2}=\res{\frac{1}{9}}$
\end{tasks}

\ex{2} \begin{tasks}(2)
  \task \res{F}: $2^3+2^4=8+16=24$
  \task \res{V}
  \task \res{F}: $(2+3)^2=25$, \ $2^2+3^2=13$
  \task \res{F}: $(3^2)^3=3^6$
  \task \res{F}: $5^0=1$
  \task \res{V}: $(6:3)^4=2^4$
\end{tasks}

\ex{3} \begin{tasks}[after-item-skip=6pt](1)
  \task $\dfrac{2^4\cdot5^4\cdot2^3}{5^3\cdot2^{10}}=\dfrac{2^7\cdot5^4}{2^{10}\cdot5^3}=\dfrac{5}{2^3}=\res{\frac{5}{8}}$
  \task $3^6\cdot3^3:3^7=3^9:3^7=3^2=\res{9}$
  \task $\dfrac{2^6\cdot3^3\cdot2^2\cdot3^2}{2^2\cdot3^4\cdot2^5}=\dfrac{2^8\cdot3^5}{2^7\cdot3^4}=2\cdot3=\res{6}$
\end{tasks}

\ex{3} $\left(\frac{4}{9}\right)^2=\left[\left(\frac{2}{3}\right)^2\right]^2=\left(\frac{2}{3}\right)^4 \ \Rightarrow\ \left(\frac{2}{3}\cdot\frac{3}{2}\right)^4=1^4=\res{1}$


% =====================================================
\seccio{10}{Potencias de exponente negativo}

\ex{1} \begin{tasks}[after-item-skip=6pt](4)
  \task $5^{-2}=\res{\frac{1}{25}}$
  \task $2^{-5}=\res{\frac{1}{32}}$
  \task $10^{-3}=\res{\frac{1}{1000}}$
  \task $\left(-3\right)^{-2}=\res{\frac{1}{9}}$
  \task $\left(-2\right)^{-3}=\res{-\frac{1}{8}}$
  \task $\left(\frac{1}{4}\right)^{-2}=\res{16}$
  \task $\left(\frac{2}{5}\right)^{-3}=\res{\frac{125}{8}}$
  \task $\left(-\frac{3}{2}\right)^{-1}=\res{-\frac{2}{3}}$
\end{tasks}

\ex{2} $-\frac{1}{2},\ \frac{1}{4},\ -\frac{1}{8},\ 1,\ 4,\ \frac{1}{2}$ \ $\Rightarrow$ \ $\res{(-2)^{-1}<(-2)^{-3}<2^{-2}<2^{-1}<2^{0}<(-2)^{2}}$

\ex{2} \begin{tasks}[after-item-skip=6pt](2)
  \task $8^{-2}=(2^3)^{-2}=\res{2^{-6}}$
  \task $\left(\frac{1}{4}\right)^{3}=(2^{-2})^{3}=\res{2^{-6}}$
  \task $\left(\frac{1}{16}\right)^{-2}=16^{2}=(2^4)^2=\res{2^{8}}$
  \task $\dfrac{4^{-1}\cdot2^{5}}{8^{2}}=\dfrac{2^{-2}\cdot2^{5}}{2^{6}}=2^{-2+5-6}=\res{2^{-3}}$
\end{tasks}

\ex{2} \begin{tasks}[after-item-skip=6pt](3)
  \task $7^{-3}\cdot7^{5}=\res{7^{2}}$
  \task $3^{4}:3^{-2}=\res{3^{6}}$
  \task $(5^{-2})^{3}=\res{5^{-6}}$
  \task $2^{-4}\cdot5^{-4}=\res{10^{-4}}$
  \task $\left(\frac{3}{4}\right)^{-2}\cdot\left(\frac{3}{4}\right)^{5}=\res{\left(\frac{3}{4}\right)^{3}}$
  \task $6^{-3}:2^{-3}=\res{3^{-3}}$
\end{tasks}

\ex{3} \begin{tasks}[after-item-skip=6pt](1)
  \task $\dfrac{9^{2}\cdot3^{-5}}{27^{-1}}=\dfrac{3^{4}\cdot3^{-5}}{3^{-3}}=3^{4-5+3}=\res{3^{2}}$
  \task $\dfrac{4^{3}\cdot6^{-2}}{3^{-4}\cdot2^{5}}=\dfrac{2^{6}\cdot2^{-2}\cdot3^{-2}}{3^{-4}\cdot2^{5}}=2^{-1}\cdot3^{2}=\res{\frac{9}{2}}$
  \task $\dfrac{10^{-2}\cdot25^{2}}{5^{3}\cdot2^{-3}}=\dfrac{2^{-2}\cdot5^{-2}\cdot5^{4}}{5^{3}\cdot2^{-3}}=2\cdot5^{-1}=\res{\frac{2}{5}}$
\end{tasks}

\ex{3} \begin{tasks}[after-item-skip=6pt](1)
  \task $2^{-1}+2^{-2}-2^{-3}=\frac{1}{2}+\frac{1}{4}-\frac{1}{8}=\frac{4}{8}+\frac{2}{8}-\frac{1}{8}=\res{\frac{5}{8}}$
  \task $\left(\frac{2}{3}\right)^{-2}-\left(\frac{1}{2}\right)^{-1}\cdot \left(-3\right)^{-1}=\frac{9}{4}-2\cdot \left(-\frac{1}{3}\right)=\frac{9}{4}+\frac{2}{3}=\frac{27}{12}+\frac{8}{12}=\res{\frac{35}{12}}$
  \task $\left[\left(\frac{1}{2}\right)^{-2}-\left(-3\right)^{0}\right]^{-1}=\left[4-1\right]^{-1}=3^{-1}=\res{\frac{1}{3}}$
  \task $\left(\frac{3}{4}-\frac{1}{2}\right)^{-2}:\left(2-\frac{4}{3}\right)^{-1}=\left(\frac{1}{4}\right)^{-2}:\left(\frac{2}{3}\right)^{-1}=16:\frac{3}{2}=\res{\frac{32}{3}}$
\end{tasks}


\end{document}
